% LSPV-3 REPORT::SGLHYB1::2026-10-02T00:29:35-07:00
% !TEX program = xelatex
\documentclass[11pt]{article}
\usepackage[margin=0.88in]{geometry}
\usepackage{amsmath,amssymb}
\begin{document}
\section*{LSPV-3 SGLHYB1: hybrid call-count/query-writing ledger for relativized SGL reconstruction}

\textbf{Global status: OPEN.} P versus NP remains open. This report sharpens only a conditional reconstruction upper-bound ledger. It proves no SAT lower bound and no hierarchy separation.

\paragraph{Target and exact prefix.}
The retained dependency-sensitive reconstruction prefix is
\[
\boxed{\exists\sigma\;\forall u\;\exists z\;\forall y\;R(x,Q,\sigma,u,z,y).}
\]
The target is not generic quantifier commutation. It is to account exactly for the cost of replacing SAT-oracle calls after the relativized Sipser--G\'acs--Lautemann (SGL) simulation, while preserving one shared size-normalized $Q$, Arthur randomness, exact YES/NO semantics, and model-matched compilation.

\paragraph{Predecessor preflight.}
The newest verified LSPV-2 artifact visible in the canonical directory remains
\texttt{2026-10-01\_1918\_LSPV-2\_SUPPORTK1.tex}, which is substantive and sentinelled. Later scheduled LSPV-2 cycles have not produced a newer verified canonical report. Existing recovery-first instructions already prohibit reconstructing missing scientific content, so this run does not invent a predecessor body and proceeds nonblocking.

\section*{CORPUS REVIEWED}
Before new work I surveyed repository search results and reviewed root \texttt{main.tex}, \texttt{AGENTS.md}, \texttt{CODEX.md}, \texttt{.codex/skills/pnp-research/SKILL.md}, current task specifications including TASK-005, current and historical MAIN material surfaced by repository search, LSPV-2 SUPPORTK1 and its WEDGE/POLAR/window/convolution predecessors, LSPV-3 DEPSWAP1/TRIVUNI1/NOGUARD1/SGLALT1/SGLAUD1/SGLCOST1/QLEN1/XOVER2, LSPV-4 WEIGHTPAD1, and the HEC/CRS-PPC/CCSPC causal branch. Repository search was also used to locate duplicate SGL/query-count formulations and stale MAIN text.

The direct primary literature checked this run was E. Viola, \emph{On Probabilistic Time versus Alternating Time}, ECCC TR05-137 (2005), and D. van Melkebeek--N. Mocelin Sdroievski, \emph{Instance-Wise Hardness and Refutation versus Derandomization for Arthur--Merlin Protocols}, Computational Complexity 34:15 (2025). Viola records
\[
\mathrm{BPTime}(t)\subseteq\Sigma_2\mathrm{Time}(t^2\log t)
\]
and
\[
\mathrm{BPTime}(t)\subseteq\Sigma_3\mathrm{Time}(t\,\mathrm{polylog}\,t),
\]
and an oracle relative to which the two-alternation simulation cannot be improved to $o(t^2)$. The positive simulations relativize. The 2025 vMS paper remains the targeted AM/reconstruction architecture, not an unrestricted deterministic lower-bound theorem.

\section*{Hybrid SGL substitution ledger}
Let the base probabilistic SAT-oracle reconstruction machine have running time $t$ in a standard explicit-query-writing oracle-Turing-machine model. Suppose per base execution it makes
\[
q(t)=t^{\delta+o(1)}
\]
SAT calls and each SAT query has length at most
\[
S_{\max}(t)=t^{b+o(1)},\qquad 0\le b\le1.
\]
Let the counterfactual exact SAT decider run in time
\[
A(s)=s^{\kappa+o(1)},\qquad \kappa\ge1.
\]

QLEN1, independently audited by LSPV-4 WEIGHTPAD1 in its declared model, gives for any explicit-query oracle computation of total time $T$ and query lengths $s_i$:
\[
\sum_i s_i\le T,
\qquad
\sum_i A(s_i)\le T S_{\max}^{\kappa-1+o(1)}.
\]

The classical SGL covering simulation evaluates only polynomially many shifted/amplified copies of the base computation. At exponent level its ordinary running time is
\[
T_{\rm SGL}(t)=t^{2+o(1)}.
\]
If the source-level implementation uses $t^{1+o(1)}$ base-machine evaluations after amplification/covering, then the inherited SAT-call count is
\[
Q_{\rm SGL}(t)=t^{1+\delta+o(1)}.
\]
This last census is a \textbf{CONDITIONAL source-level parameterization} until the exact amplified implementation used in the project is transcribed and independently audited. It is not restored as the old SGLAUD1 claim.

Two independent substitution upper bounds are then available:
\[
\sum_i A(s_i)
\le Q_{\rm SGL}S_{\max}^{\kappa+o(1)}
=t^{1+\delta+b\kappa+o(1)}
\]
and
\[
\sum_i A(s_i)
\le T_{\rm SGL}S_{\max}^{\kappa-1+o(1)}
=t^{2+b(\kappa-1)+o(1)}.
\]
Hence, under the stated source-level census,
\[
\boxed{
c_{\rm hyb}(\kappa,b,\delta)
=
\min\{1+\delta+b\kappa,\;2+b(\kappa-1)\}
=
1+b\kappa+\min\{\delta,1-b\}.
}
\]
The algebraic equality is exact. The specialized reconstruction application is \textbf{CONDITIONAL} because $b,\delta$ and the exact SGL census remain source/model-matching obligations.

\paragraph{Regime split.}
The call-count bound wins exactly when
\[
\delta<1-b,
\]
the query-writing amortization wins exactly when
\[
\delta>1-b,
\]
and they tie at $\delta=1-b$. Thus a small maximum query-size exponent can make a genuine base-query census useful; when $b=1$, query-writing amortization is never worse for any $\delta\ge0$ and gives $c_{\rm hyb}=\kappa+1$.

If two later model-matched SAT-compilation/elimination stages contribute exponents $b_1,b_2$, the resulting candidate SGL ledger is
\[
\boxed{
\alpha_{\rm SGL}^{\rm hyb}
\le
\alpha_0 b_1b_2\kappa^2
\left(1+b\kappa+\min\{\delta,1-b\}\right)+o(1).
}
\]
This is \textbf{CONDITIONAL upper-bound accounting}. It is not a lower bound and does not establish optimality.

\paragraph{Comparison with the conservative route.}
The dependency-sensitive conservative ledger remains
\[
\alpha_{\rm V}
\le
\alpha_0(b_0b_1b_2b_3)\kappa^4+o(1).
\]
The correct current envelope is the minimum of independently justified routes after parameter identification. XOVER2's compiler-sensitive phase diagram remains algebraically correct for its declared simplified SGL term $\kappa+1$, but it is not the most general ledger when $S_{\max}=t^{b+o(1)}$ with $b<1$ and an independently certified base-call exponent $\delta$ is available.

\section*{Exact reconstruction semantics and charges}
For reconstruction answer vector $a$ and query family $J$,
\[
\operatorname{ExactAns}(a;Q,i,r)
=
\exists Y^+\forall Y^-
\bigwedge_{j\in J}
\left[
(a_j\Rightarrow\Phi_j(Y_j^+))
\wedge
(\neg a_j\Rightarrow\neg\Phi_j(Y_j^-))
\right].
\]
YES witnesses and NO universal assignments remain distinct. Generic canonical binding is not free. A construction-specific collision/local-algebraic binding theorem would be relevant only if it beats the charged $A\circ D_j$ baseline.

The full resource ledger remains
\[
\tau_Q+\tau_{\rm norm}+\tau_{\rm census}
+\tau_{\rm amp}+\tau_{\rm idx}
+\tau_{\rm compile}+\tau_{\rm rec}
+\tau_{\rm write}+\tau_{\rm transport}.
\]
TransportCert remains accounting only. The hybrid bound uses both $\tau_{\rm census}$ and $\tau_{\rm write}$; neither is free.

\section*{Hostile controls and barrier audit}
SAME\_REALIZATION fixes the normalized $Q$, Arthur randomness, legal reconstruction family, compiler, oracle-machine interface, and downstream SAT instances. A unit-cost pointer oracle naming long formulas without explicit writing is a model mismatch for QLEN1 and must be separately compiled and charged. DEAD/COPY/DELAY do not increase semantic burden. INTERP/RECOMP, compact alternate representations, memoization, padding, Horn/2-SAT, affine/Tseitin, bounded treewidth, independent products, semantic oracles, and algebraic oracles remain controls.

The SGL simulation, call-count census, and query-writing amortization all relativize. They therefore cannot themselves supply the required nonrelativizing hinge. They are also compatible with algebraic oracle access and do not cross the algebrization barrier. TC4 remains OPEN because an arbitrary exact SAT solver need not materialize the preferred reconstruction or algebraic representation.

\section*{PQ / cross-domain classification}
\textbf{EXACT OBJECT/THEOREM TRANSFER:} fixed-code Schur products; typed shortening/puncturing; quadratic-relation kernels; Boolean polarization/exterior identities; cyclic convolution in its declared model; Gaussian elimination; explicit query-writing amortization.

\textbf{ANALOGY/CANDIDATE:} Hilbert/syzygy and graded-Betti diagnostics; minimum-rank relation spectra; DQBF dependency machinery; construction-specific collision/special-soundness extraction; LWE-style external black-box extraction.

\textbf{FAILED TRANSFER:} cryptographic hardness assumptions, random oracles, VLSI area-time, generic control/observability, distributed cuts, topology, and proof-system-specific lower bounds as unconditional unrestricted SAT-runtime lower bounds.

\section*{DUPLICATION / CONTRADICTION AUDIT}
SGLAUD1's earlier amplified constant-query justification remains \textbf{FALSIFIED/RETIRED}; this report does not revive it. SGLCOST1 correctly forced sequential SAT-oracle substitution. QLEN1 supplied the sharper query-writing amortization. XOVER2 compared the $b=1$ specialization against the conservative route. The new content is the two-resource minimum that retains both an independently certifiable call exponent $\delta$ and a potentially sublinear maximum query-size exponent $b$.

No contradiction was found with SUPPORTK1 or the upstream Schur/exterior program. The causal HEC/CRS-PPC/CCSPC branch supplies no missing static-SAT transfer and remains a hostile-control laboratory.

\section*{WHAT CHANGED RELATIVE TO ROOT MAIN}
Root \texttt{main.tex} is stale in the LSPV-3 region: it still presents SGLAUD1's constant-query/amplification rationale and asks only for $(q(t),S_{\max}(t))$ before comparing ledgers. The corrected state is: SGLAUD1's rationale is retired; QLEN1's local explicit-query inequality is independently audited; the reconstruction ledger should retain the hybrid pair $(b,\delta)$ and use
\[
1+b\kappa+\min\{\delta,1-b\}
\]
only after the exact amplified SGL census is source-certified.

\section*{ADOPTED / REJECTED / RETIRED PRIOR COMPONENTS}
\textbf{Adopted:} DEPSWAP1 dependency guard; TRIVUNI1 $A\circ D_j$ baseline; NOGUARD1 certificate guard; SGLALT1 architecture; SGLCOST1 sequential-substitution correction; QLEN1 explicit-query amortization; XOVER2 envelope discipline.

\textbf{Rejected:} free canonical binding; free qualitative PH collapse as an exponent-preserving compiler; unit-cost long-query pointers without model transfer; any claim that algebraic structure itself forces runtime.

\textbf{Retired:} SGLAUD1's unverified amplified constant-query derivation and generic quantifier swapping.

\section*{Strict status and next target}
\textbf{KNOWN THEOREM:} classical SGL $\Sigma_2$ simulation and Viola's relativized quadratic obstruction; vMS targeted AM/reconstruction architecture.

\textbf{CONDITIONAL:} the hybrid SGL reconstruction ledger and its $Q_{\rm SGL}=t^{1+\delta+o(1)}$ source-level census pending exact transcription and independent red-team review.

\textbf{OPEN:} exact vMS/SGL parameters $b,\delta,b_1,b_2$; all SAME\_REALIZATION charge terms; TC4; nonrelativizing/nonalgebrizing hinge; final contradiction.

The single next target is to source-transcribe one actual amplified vMS reconstruction execution in the explicit-query SAT-oracle model and certify
\[
\boxed{b,\;\delta,\;b_1,\;b_2,\;\tau_{\rm amp},\;\tau_{\rm write}.}
\]
Then select the smaller of the call-count and query-writing substitution bounds before comparing against the conservative route.

\begin{thebibliography}{2}
\bibitem{Viola05} E. Viola, \emph{On Probabilistic Time versus Alternating Time}, ECCC TR05-137, 2005.
\bibitem{vMS25} D. van Melkebeek and N. Mocelin Sdroievski, \emph{Instance-Wise Hardness and Refutation versus Derandomization for Arthur--Merlin Protocols}, Computational Complexity 34:15, 2025.
\end{thebibliography}
\end{document}
% END_OF_VERIFIED_REPORT::LSPV-3::SGLHYB1
